\subsection{Convex and symplectic notation}\label{subsec:preliminaries-polytopes}

We write \(\mathbb{R}^4\) with coordinates \((q_1,q_2,p_1,p_2)\). The standard
complex structure is
\[
  J_0 =
  \begin{pmatrix}
    0 & -I_2 \\
    I_2 & 0
  \end{pmatrix},
\]
and the standard symplectic form is
\[
  \omega_0(u,v) = \langle J_0u,v\rangle
  = \sum_{i=1}^2 (u_{q_i}v_{p_i}-u_{p_i}v_{q_i}).
\]
We will use \(J_0^2=-I\), \(J_0^{\mathsf T}=-J_0\), and
\(\omega_0(J_0u,J_0v)=\omega_0(u,v)\) without further comment.
We use the standard primitive
\[
  \lambda_0 = \frac{1}{2}\sum_{i=1}^2(q_i\,dp_i-p_i\,dq_i).
\]
Thus \(d\lambda_0=\omega_0\) and
\(\lambda_0|_x(v)=\frac12\langle -J_0v,x\rangle\).
Thus, for an absolutely continuous closed curve
\(\gamma\colon[0,T]\to\mathbb{R}^4\), its action is
\begin{equation}\label{eq:preliminaries-action}
  A(\gamma) \coloneqq \int_\gamma \lambda_0
  = \frac{1}{2}\int_0^T
    \langle -J_0\dot{\gamma}(t),\gamma(t)\rangle\,dt.
\end{equation}
This action is invariant under orientation-preserving reparametrization.

\begin{lemma}[Symplectic shoelace formula]\label{lem:preliminaries-shoelace}
Let \(z\colon[0,T]\to\mathbb{R}^4\) be a closed piecewise affine curve. Suppose
that \(z\) has constant velocity \(w_k\) on consecutive intervals of length
\(\ell_k\), for \(k=1,\ldots,m\). Then
\begin{equation}\label{eq:preliminaries-shoelace}
  2A(z)
  =
  \sum_{1\leq j<k\leq m}
    \ell_j\ell_k\,\omega_0(w_j,w_k).
\end{equation}
In particular, the action of a closed pure-word Reeb orbit depends only on the
cyclic velocity word and the dwell times, not on the base point.
\end{lemma}

\begin{proof}
Write \(t_k=\sum_{j<k}\ell_j\). On the \(k\)-th interval,
\[
  z(t)=z(0)+\sum_{j<k}\ell_jw_j+(t-t_k)w_k.
\]
The contribution of \(z(0)\) to
\(\int_0^T\langle -J_0\dot z,z\rangle\,dt\) vanishes because \(z\) is closed,
so \(\sum_k\ell_kw_k=0\). The self-term
\(\langle -J_0w_k,w_k\rangle\) vanishes by skew-symmetry. Hence
\[
  2A(z)
  =
  \sum_k
    \ell_k\left\langle -J_0w_k,\sum_{j<k}\ell_jw_j\right\rangle
  =
  \sum_{j<k}\ell_j\ell_k\,\omega_0(w_j,w_k),
\]
using \(\omega_0(u,v)=\langle J_0u,v\rangle\).
\end{proof}

